\section{Composición de la holonomía \(G_9\) con la carta \(T_1\) y certificado de precisión arbitraria para \(\alpha_{\rm HMT}\)}
\label{sec:alpha-t1-10000-rev9}

La prolongación T1 aplica la recurrencia de acarreo a las tríadas
\(P_k,E_k,\Phi_k\) producidas previamente por G9 y a la coordenada
dodecafásica declarada \(K\):
\begin{align}
 s_k&=P_k+E_k-\Phi_k-K_{k\bmod12}+c_{k+1},\\
 A_k&=s_k\bmod1000,\\
 c_k&=\left\lfloor\frac{s_k}{1000}\right\rfloor.
 \label{eq:composicion-g9-t1-alpha-10000-rev9}
\end{align}
El dominio de entrada contiene las salidas ternarias generadas por G9; no
contiene una expansión decimal de \(\alpha\), un valor de
\(\alpha^{-1}\), CODATA ni una condición metrológica objetivo.

Los cilindros de \(3501\) bloques fijan \(10020\) cifras de guarda en cada
canal, equivalentes a \(3340\) tríadas. El cálculo propaga, de derecha a
izquierda, todas las condiciones terminales posibles del conjunto estable
\[
 \mathcal C=\{-2,-1,0,1,2\}.
\]
Las cinco ramas coalescen en un prefijo común de \(3339\) tríadas, es decir,
\(10017\) cifras. En consecuencia, las primeras diez mil cifras de
\(\alpha_{\rm HMT}\) son independientes del acarreo situado más allá de la
zona de guarda. La salida publicada no impone \(c_N=0\).

La independencia causal puede comprobarse después de la construcción: las
primeras tres mil cifras coinciden con la expansión decimal obtenida por el
lector externo, pero ninguna de ellas interviene en
\eqref{eq:composicion-g9-t1-alpha-10000-rev9}. La expansión de diez mil
cifras y los datos finitos que permiten repetir el cálculo se muestran en el
apéndice~\ref{app:desarrollos-10000-rev9}.
